\chapter{Теоретические сведения}
\label{ch:theory}

\section{Ansible — система управления конфигурациями}

\textbf{Ansible} — свободная платформа автоматизации, разработанная компанией
Red~Hat. Ключевые особенности:
\begin{itemize}
  \item \textbf{Agentless} — не требует установки агентов на управляемые узлы;
        подключение осуществляется по \textbf{SSH};
  \item \textbf{Idempotent} — повторный запуск задачи не изменяет уже корректно
        настроенную систему;
  \item \textbf{YAML-синтаксис} — плейбуки и роли описываются в формате,
        удобном для чтения человеком;
  \item \textbf{Модульность} — более 3000 встроенных модулей для управления
        пакетами, сервисами, файлами, сетью и т.д.
\end{itemize}

\section{Архитектура Ansible}

Основные компоненты:
\begin{itemize}
  \item \textbf{Control node} — машина, на которой установлен Ansible
        и с которой запускаются плейбуки. В работе — \texttt{gateway}.
  \item \textbf{Managed nodes} — узлы, которыми управляет Ansible.
        Требуется только \texttt{python3} и доступ по SSH.
  \item \textbf{Inventory} — файл (\texttt{/etc/ansible/hosts}),
        перечисляющий управляемые узлы и группы.
  \item \textbf{Playbook} — YAML-файл с набором задач (\textit{tasks}),
        применяемых к группе хостов.
  \item \textbf{Module} — атомарная единица работы Ansible
        (\texttt{file}, \texttt{copy}, \texttt{set\_fact} и др.).
\end{itemize}

\section{SSH-аутентификация и ключи}

Ansible подключается к managed nodes по SSH. Рекомендуемый способ — ключ
\texttt{ed25519}. Алгоритм EdDSA (Edwards-curve Digital Signature Algorithm)
предпочтительнее RSA по двум причинам:
\begin{enumerate}
  \item \textbf{Скорость} — операции подписи и верификации значительно быстрее;
  \item \textbf{Безопасность} — устойчивость к атакам на реализацию
        (нет зависимости от генератора случайных чисел при подписи).
\end{enumerate}

Генерация ключа:
\begin{lstlisting}
ssh-keygen -t ed25519 -a 100
\end{lstlisting}

Распространение публичного ключа на клиент:
\begin{lstlisting}
ssh-copy-id username@ip_or_hostname
\end{lstlisting}

\section{Inventory-файл}

Файл \texttt{/etc/ansible/hosts} определяет, какими узлами управляет Ansible.
Пример структуры:
\begin{lstlisting}
[clients]
desktop1 ansible_host=192.168.29.10 ansible_user=username
wordpress ansible_host=192.168.29.6  ansible_user=username
\end{lstlisting}

Здесь \texttt{clients} — имя группы; \texttt{ansible\_host} — IP или
DNS-имя; \texttt{ansible\_user} — пользователь на управляемом узле.

\section{Структура плейбука и модули gather\_facts}

\textbf{Плейбук} начинается с заголовка \texttt{-~name} и содержит директивы
\texttt{hosts} (группа из inventory), \texttt{gather\_facts}~(автосбор
системных параметров) и список задач \texttt{tasks}.

Директивы делегирования:
\begin{itemize}
  \item \texttt{delegate\_to: localhost} — задача выполняется
        не на клиенте, а на control node;
  \item \texttt{run\_once: true} — задача выполняется только один раз
        (например, создание папки на сервере).
\end{itemize}

Переменные \texttt{ansible\_facts}, используемые в работе:
\begin{itemize}
  \item \texttt{ansible\_hostname} — имя хоста клиента;
  \item \texttt{ansible\_default\_ipv4.address} — основной IP-адрес;
  \item \texttt{ansible\_distribution}, \texttt{ansible\_distribution\_version}
        — название и версия ОС;
  \item \texttt{ansible\_mounts} — список примонтированных файловых систем
        со статистикой (в том числе \texttt{size\_available} для~\texttt{/}).
\end{itemize}

\endinput
